% Mandatory Data in Brief Specifications Table.
% Page-breakable longtable: the table is taller than one page, so a fixed float
% (table*) overflows the type block. longtable breaks cleanly across pages.
% All entries MUST match paper/_shared/CANONICAL_FACTS.md and the published CSVs.
\begingroup
\renewcommand{\arraystretch}{1.3}
\begin{longtable}{@{}>{\raggedright\arraybackslash}p{0.26\textwidth}p{0.66\textwidth}@{}}
\caption{Specifications Table.}
\label{tab:specifications}\\
\toprule
Subject & Aerospace Engineering \\
\midrule
Specific subject area & Supersonic rocket aerodynamics and trajectory simulation; apogee
prediction and model validation for sounding rockets and high-power amateur rocketry. \\
\midrule
Type of data & Table (comma-separated values, CSV); plain-text provenance documentation
(Markdown). The reusable observational core is the externally measured flight apogees
(raw) together with the independently published RASAero~II reference predictions; the
this-work apogee columns are reproducible simulated predictions supplied as paired
benchmark metadata (processed), not the primary observed datum. \\
\midrule
How the data were acquired &
Externally measured flight ground truth was compiled from published instrumented flight
records: barometric altimeter, GPS, optical (theodolite) track, integrated accelerometer,
and radar/beacon tracking. Paired commercial-simulator reference apogees are the recorded
RASAero~II~\cite{rogers2015} altitude-comparison predictions published by Rogers
Aeroscience~\cite{rogers_rasaero_alt}. The OpenRocket-Plus (\texttt{apogee\_thiswork\_ft})
column was generated by running the archived simulator release on each reconstructed
vehicle; predicted apogees are read directly from the simulation output. The reusable
observational core of the corpus is therefore the externally measured flight apogees plus
the independently published RASAero~II predictions (both observation/external); the
this-work apogee column is a reproducible simulated prediction supplied as paired benchmark
metadata, not the primary observed datum. \\
\midrule
Data format &
Raw and analysed.
Two CSV tables (\texttt{flight\_comparison.csv}, \texttt{sounding\_rockets\_exploratory.csv})
and one Markdown provenance/reproducibility statement (\texttt{reproducibility.md}). \\
\midrule
Description of data collection &
The headline corpus comprises 25 externally curated flights spanning Mach~0.54--4.33
(23 single-stage high-power flights plus two two-stage vehicles, the AeroPac~104K Two-Stage
at Mach~3.04 and the MESOS~293K closure at Mach~4.33), each pairing a
measured apogee with the corresponding RASAero~II reference prediction. A separate
exploratory set of $\sim$20 historical high-Mach sounding-rocket flights (e.g., Nike-Apache,
Nike-Deacon, Nike-Cajun, Arcas, Black~Brant~V, HEROS~3) is reported in full, with primary
flight-report citations per row. The headline corpus was selected externally (it is Rogers'
public RASAero~II comparison set), not curated by outcome. \\
\midrule
Data source location &
Flights are drawn from published sources documented per record in the dataset. Headline
measured and RASAero~II values: Rogers Aeroscience RASAero~II altitude-comparison
set~\cite{rogers_rasaero_alt}. Exploratory ground truth: original NACA/NASA/DTIC flight
reports (cited per row, e.g.,~\cite{heitkotter1956,nasa_x721_66_568,dtic_ad0733141}).
Geographic launch sites are recorded per flight (e.g., Black Rock Desert, NV, USA;
Wallops Island, VA, USA; White Sands, NM, USA; Esrange, Sweden; Churchill, MB, Canada). \\
\midrule
Data accessibility &
Repository name: Rocket Flight Database. \newline
Data identification number (DOI): 10.5281/zenodo.20531977 (version of record at submission). \newline
Direct URL to data: \url{https://doi.org/10.5281/zenodo.20531977}. \newline
Source-code/development mirror: \url{https://github.com/AidanSYu/rocket-flight-database}. \newline
License: Creative Commons Attribution 4.0 International (CC-BY-4.0)~\cite{rfd_zenodo}. \\
\midrule
Related research article &
A. Yu, ``Integrated open-source supersonic aerodynamics for sounding-rocket and high-power
apogee prediction: a 25-flight validation and architecture,'' \textit{Journal of Spacecraft
and Rockets} (companion research article, under review)~\cite{yu2026jsr}. The present dataset is the
validation corpus described and analysed in that article; this Data in Brief article describes
the data only and draws no new scientific conclusions. \\
\bottomrule
\end{longtable}
\endgroup
